%Sistema Internacional De Unidades:
\NC{\n}[2][]{\, \txnor{#2}^{#1}}
\NCX{\nfr}[4][1=,3=]{\, \fr{\txnor{#2}^{#1}}{\txnor{#4}^{#3}}}
\NCX{\nfrr}[4][1=,3=]{\, \frr{\txnor{#2}^{#1}}{\txnor{#4}^{#3}}}
\NC{\nfrac}[2]{\, \fr{#1}{#2}}
\NC{\nfrrac}[2]{\, \frr{#1}{#2}}
\NC{\npor}{\tresl \por \cincol}

%\NCX{\p}[2][1=,2=]{\partial_{#1}^{#2}} %Derivada Parcial.
%Submúltiplo En Micrómetros:
\NC{\micro}{\tx{\textmu}}

%Operadores:
\NC{\op}{\hat}
%Conmutador:
\NC{\conm}[2]{\lb{#1,\dosl #2}\rb}

%Brakets:
\NC{\bra}[1]{\la{#1}\rv}
\NC{\ket}[1]{\lv{#1}\ra}
\NC{\bret}[2]{\la #1 \middle| #2 \ra}
\NC{\braket}[3]{\la #1 \middle| #2 \middle| #3 \ra}

%Circuitos Eléctricos:

%Dimensiones De Componentes Eléctricos:
\ctikzset{bipoles/resistor/height=0.25, bipoles/resistor/width=0.6} %Resistencia.
\ctikzset{bipoles/capacitor/height=0.525, bipoles/capacitor/width=0.2} %Capacitor.
\ctikzset{bipoles/cuteinductor/width/.initial=0.7, bipoles/cuteinductor/coils/.initial=7} %Inductancia.
\ctikzset{bipoles/diode/height=0.35, bipoles/diode/width=0.3} %Diodo.
\ctikzset{bipoles/lamp/height=0.45, bipoles/lamp/width=0.45} %Lámpara.
\ctikzset{bipoles/thermistor/height=0.3, bipoles/thermistor/width=0.7} %Termistor.
\ctikzset{amplifiers/scale=0.2, amp symbol font={\footnotesize}}
%Letras Barra:
\NC{\hj}{\hslash} %Constante De Planck Reducida.
\NC{\lj}{{\mkern0.75mu\mathchar '26\mkern -9.75mu\lambda}} %Longitud De Onda Angular.

%Constantes Físicas:
\NC{\kr}{k} %Constante Del Resorte.
\NC{\kgu}{G} %Constante De Gravitación Universal.
\NC{\kc}{k_{\tx{e}}} %Constante Electrostática (Constante De Coulomb).
\NC{\kpev}{\eps_0} %Constante De Permitividad Eléctrica Del Vacío.
\NC{\kpmv}{\mi_0} %Constante De Permeabilidad Magnética Del Vacío.
\NC{\kse}{\ji_{\tx{e}}} %Constante De Susceptibilidad Eléctrica.
\NC{\ksm}{\ji_{\tx{m}}} %Constante De Susceptibilidad Magnética.
\NC{\kb}{k_{\tx{B}}} %Constante De Boltzmann.
\NC{\ka}{N_{\tx{A}}} %Constante De Avogadro.
\NC{\kgi}{R} %Constante De Los Gases Ideales.
\NC{\kp}{h} %Constante De Planck.
\NC{\kpr}{\hj} %Constante De Planck Reducida.
\NC{\ksb}{\sigma} %Constante De Stefan-Boltzmann.